%!TEX root = ../../main.tex
\section{구간 변화와 전략적 가치 개선의 정의}
\label{sec:value-svi}

\begin{definition}[구간 변화와 SVI]
예측 시점 이전 상태 $\xpre$와 결과 시점 상태 $\xpost$(식 \eqref{eq:endpoint}의 $\eh$, $h = 90$)에 대해
\begin{equation}
  \dV = \Vhat(\xpost) - \Vhat(\xpre), \qquad \Ysvi = \ind[\dV > 0].
  \label{eq:svi}
\end{equation}
$\Ysvi = 0$은 감소와 정확히 변하지 않은 경우($\dV = 0$)를 모두 포함한다. 전·후 점수가 유한하지 않은 교전은 0으로 부호화하지 않고 라벨 전에 제외한다. 15.14 out-of-fold 학습 라벨에는 정확한 0이 하나 있다.
\end{definition}

$\Ysvi = 1$은 동결 평가기 아래에서 \emph{추정된} 블루 승률이 교전 구간 동안 올랐다는 뜻이다. 변화량은 $100\,\dV$ 퍼센트포인트(예: $0.70 \to 0.60$은 $-10$ 퍼센트포인트)이고, 같은 창 안에 다른 상태 갱신이 들어올 수 있으므로 $\dV$는 구간의 변화이며 교전 하나의 인과 기여와 구분된다(성분 분해는 제 \ref{sec:value-shold} 절). 주 예측 목표는 \emph{방향} $\Ysvi$이며 $\dV$의 크기 회귀는 범위 밖이다.

균형 구간 $\Bforty$는 동결 fit85 $\ppre$로 다시 계산한 $\ppre \in [0.40, 0.60]$인 교전(0.5를 중심으로 폭 20퍼센트포인트)이며, 15.16 TEST에서 T는 \numTBforty{} 교전(\numTBfortyMatches{} 경기, T의 \shareTBforty\%), S는 \numSBforty{} 교전(\numSBfortyMatches{} 경기, S의 \shareSBforty\%)이다. 전체 코호트와 나란히 보고하는 조건부 평가이다.

교환가치 라벨과의 관계는 같은 사례에서 두 라벨을 계산한 제 \ref{sec:value-samecase} 절에서 다룬다.
